\documentclass[12pt,a4paper]{report}

% -------------------------------------------------------------
% UNIVERSAL FONT & COMPILER SETUP (XeLaTeX / LuaLaTeX / pdfLaTeX)
% -------------------------------------------------------------
\usepackage{iftex}
\ifxetex
  \usepackage{fontspec}
  \setmainfont{Liberation Serif}
\else\ifluatex
  \usepackage{fontspec}
  \setmainfont{Liberation Serif}
\else
  \usepackage[utf8]{inputenc}
  \usepackage[T1]{fontenc}
  \usepackage{mathptmx}
\fi\fi

% -------------------------------------------------------------
% PAGE GEOMETRY & SPACING (Matching College Template)
% -------------------------------------------------------------
\usepackage[left=1.25in,right=1in,top=1in,bottom=1in]{geometry}
\usepackage{setspace}
\setstretch{1.25}
\setlength{\parindent}{1.5em}
\setlength{\parskip}{0.3em}
\raggedbottom

% Float placement optimization to eliminate empty whitespace
\renewcommand{\topfraction}{0.9}
\renewcommand{\bottomfraction}{0.8}
\renewcommand{\textfraction}{0.1}
\renewcommand{\floatpagefraction}{0.75}
\setlength{\textfloatsep}{12pt plus 2pt minus 2pt}
\setlength{\floatsep}{10pt plus 2pt minus 2pt}
\setlength{\intextsep}{10pt plus 2pt minus 2pt}

% -------------------------------------------------------------
% MATHEMATICS & ALGORITHMS
% -------------------------------------------------------------
\usepackage{amsmath,amssymb,amsfonts}
\usepackage{bm}

% -------------------------------------------------------------
% GRAPHICS, TABLES & CAPTION STYLING
% -------------------------------------------------------------
\usepackage{graphicx}
\graphicspath{{figures/}{report/figures/}}
\usepackage{booktabs}
\usepackage{longtable}
\usepackage{tabularx}
\usepackage{array}
\usepackage{caption}
\usepackage{subcaption}
\captionsetup{font={small,rm},labelfont=bf,skip=5pt}

% -------------------------------------------------------------
% HEADERS & FOOTERS (Matching College Template)
% -------------------------------------------------------------
\usepackage{fancyhdr}
\pagestyle{fancy}
\fancyhf{}
\fancyhead[R]{\small\nouppercase{\leftmark}}
\fancyfoot[C]{\thepage}
\renewcommand{\headrulewidth}{0.4pt}
\renewcommand{\footrulewidth}{0pt}
\setlength{\headheight}{15pt}

\fancypagestyle{plain}{
  \fancyhf{}
  \fancyfoot[C]{\thepage}
  \renewcommand{\headrulewidth}{0pt}
}

% -------------------------------------------------------------
% CHAPTER AND SECTION HEADING FORMATTING (Matching College Template)
% -------------------------------------------------------------
\usepackage{titlesec}
\titleformat{\chapter}[display]
  {\normalfont\Large\bfseries\centering}
  {\MakeUppercase{\chaptertitlename}\ \thechapter}
  {10pt}
  {\Large\bfseries\MakeUppercase}
\titlespacing*{\chapter}{0pt}{-10pt}{25pt}

\titleformat{\section}
  {\normalfont\large\bfseries}{\thesection}{1em}{}
\titlespacing*{\section}{0pt}{14pt}{6pt}

\titleformat{\subsection}
  {\normalfont\normalsize\bfseries}{\thesubsection}{1em}{}
\titlespacing*{\subsection}{0pt}{10pt}{4pt}

% -------------------------------------------------------------
% HYPERREF SETUP
% -------------------------------------------------------------
\usepackage{xcolor}
\definecolor{navyblue}{RGB}{20,50,110}
\usepackage[
  colorlinks=true,
  linkcolor=navyblue,
  citecolor=navyblue,
  urlcolor=navyblue,
  pdfstartview=FitH
]{hyperref}

% -------------------------------------------------------------
% DOCUMENT METADATA DEFINITIONS
% -------------------------------------------------------------
\newcommand{\projecttitle}{COURSE PROJECT TITLE IN ALL CAPS}
\newcommand{\shorttitle}{Short Project Title}
\newcommand{\coursecode}{PBXXXX00}
\newcommand{\coursename}{COURSE NAME}
\newcommand{\degree}{Bachelor of Technology}
\newcommand{\branch}{Artificial Intelligence and Data Science}
\newcommand{\department}{Department of Artificial Intelligence and Data Science}
\newcommand{\institution}{Muthoot Institute of Technology and Science (MITS)}
\newcommand{\affiliation}{(Autonomous, Affiliated to APJ Abdul Kalam Technological University)}
\newcommand{\location}{Varikoli, Ernakulam, Kerala}
\newcommand{\monthyear}{OCTOBER 2026}
\newcommand{\academicyear}{2026--27}
\newcommand{\facultyincharge}{Faculty Name}

\begin{document}

% =============================================================
% PAGE 1: TITLE PAGE (Matching College Template)
% =============================================================
\begin{titlepage}
\thispagestyle{empty}
\centering
\setstretch{1.2}

\vspace*{0.3cm}
{\Large\bfseries \projecttitle \par}
\vspace{1.2cm}

{\large A COURSE PROJECT REPORT \par}
\vspace{0.4cm}
{\normalsize submitted by \par}
\vspace{0.4cm}
{\large\bfseries STUDENT ONE (MUT24XX000) \par}
{\large\bfseries STUDENT TWO (MUT24XX000) \par}
{\large\bfseries STUDENT THREE (MUT24XX000) \par}
{\large\bfseries STUDENT FOUR (MUT24XX000) \par}
\vspace{1.0cm}

{\normalsize as part of the course \par}
\vspace{0.3cm}
{\large\bfseries \coursecode\ -- \coursename \par}
\vspace{1.0cm}

{\normalsize to \par}
\vspace{0.3cm}
{\normalsize The APJ Abdul Kalam Technological University in partial fulfilment of the\\
requirements for the award of the Degree of \par}
\vspace{0.3cm}
{\large\bfseries \degree \par}
{\normalsize in \par}
{\large\bfseries \branch \par}
\vfill
\includegraphics[width=2.0in]{figures/college_logo.png} \par
\vspace{0.4cm}
{\normalsize \department \par}
\vspace{0.2cm}
{\large\bfseries \institution \par}
{\normalsize \location \par}
\vspace{0.4cm}
{\normalsize \monthyear \par}
\end{titlepage}

% =============================================================
% PAGE 2: CERTIFICATE (Matching College Template)
% =============================================================
\newpage
\thispagestyle{empty}
\setstretch{1.3}

\begin{center}
    \vspace*{1.0cm}
    {\Large\bfseries CERTIFICATE \par}
    \vspace{1.2cm}
\end{center}

\noindent This is to certify that the report entitled \textbf{``\projecttitle''}, submitted by \textbf{STUDENT ONE (MUT24XX000)}, \textbf{STUDENT TWO (MUT24XX000)}, \textbf{STUDENT THREE (MUT24XX000)}, and \textbf{STUDENT FOUR (MUT24XX000)} to \textbf{\institution} for the award of the degree of \textbf{\degree} in \textbf{\branch}, is a bonafide record of the course project work carried out by them, as part of the course \textbf{\coursecode\ -- \coursename}, under my supervision and guidance.

\vspace{3.2cm}

\noindent
\begin{tabularx}{\textwidth}{@{}X r@{}}
\textbf{Place:} Varikoli & \textbf{Faculty-in-Charge} \\
\textbf{Date:} 22/10/2026 & \department
\end{tabularx}

% =============================================================
% PRELIMINARY PAGES: ABSTRACT, CO-PO MATRIX & SDG (Compact)
% =============================================================
\newpage
\pagenumbering{roman}
\setcounter{page}{2}
\setstretch{1.15}

\begin{center}
    {\large\bfseries ABSTRACT \par}
    \vspace{0.2cm}
\end{center}
\addcontentsline{toc}{chapter}{Abstract}

\noindent [Abstract text goes here -- a concise, self-contained summary of the problem background, mathematical/computational methods, dataset scope, key simulation results, and practical engineering significance.]

\vspace{0.5cm}
\begin{center}
    {\large\bfseries CO -- PO / PSO ARTICULATION MATRIX \par}
    \vspace{0.2cm}
\end{center}
\addcontentsline{toc}{chapter}{CO -- PO / PSO Articulation Matrix}

\noindent Course Outcomes (COs) addressed by this project:
\begin{itemize}
  \setlength{\itemsep}{0pt}
  \setlength{\parskip}{0pt}
  \item \textbf{CO1:} Explain fundamental graph concepts and algorithms to analyse and solve problems.
  \item \textbf{CO2:} Illustrate matching and covering problems in graphs using algorithms and theorems.
  \item \textbf{CO3:} Apply connectivity and paths including Network Flow and Max-Flow Min-Cut Theorem.
  \item \textbf{CO4:} Illustrate Vertex Coloring, Upper Bounds, Brooks' Theorem, and Color-Critical Graphs.
  \item \textbf{CO5:} Apply concepts of planar graphs, Kuratowski's and Wagner's theorems, TSP, and dominating sets.
\end{itemize}

\begin{table}[htbp]
\centering
\caption{Course Outcomes to Program Outcomes Articulation Matrix}
\label{tab:co_po_matrix_template}
{\small\setlength{\tabcolsep}{2.5pt}
\begin{tabularx}{\textwidth}{l*{12}{>{\centering\arraybackslash}X}}
\toprule
\textbf{CO} & \textbf{PO1} & \textbf{PO2} & \textbf{PO3} & \textbf{PO4} & \textbf{PO5} & \textbf{PO6} & \textbf{PO7} & \textbf{PO8} & \textbf{PO9} & \textbf{PO10} & \textbf{PO11} & \textbf{PO12} \\
\midrule
CO1 & 3 & 3 & 3 & 3 & - & - & - & - & - & - & - & - \\
CO2 & 3 & 3 & 3 & - & - & - & - & - & - & - & - & - \\
CO3 & 3 & 3 & 3 & 3 & - & - & - & - & - & - & - & - \\
CO4 & 3 & 3 & 3 & 3 & - & - & - & - & - & - & - & - \\
CO5 & 3 & 3 & 3 & - & - & - & - & - & - & - & - & - \\
\bottomrule
\end{tabularx}}
\end{table}

\vspace{0.4cm}
\begin{center}
    {\large\bfseries SUSTAINABLE DEVELOPMENT GOALS (SDG) ADDRESSED \par}
    \vspace{0.2cm}
\end{center}
\addcontentsline{toc}{chapter}{Sustainable Development Goals (SDG) Addressed}

\begin{table}[htbp]
\centering
\caption{Sustainable Development Goals Alignment}
\label{tab:sdg_template}
{\small\setlength{\tabcolsep}{4pt}
\begin{tabularx}{\textwidth}{>{\raggedright\arraybackslash}p{2.6cm}>{\raggedright\arraybackslash}p{3.8cm}X}
\toprule
\textbf{Goal} & \textbf{Target} & \textbf{Justification} \\
\midrule
\textbf{SDG 11:}\newline Sustainable Cities & 
\textbf{Target 11.2:}\newline Safe, sustainable transport systems. & 
[Project contribution to sustainable mobility and infrastructure optimization.] \\
\addlinespace
\textbf{SDG 9:}\newline Industry \& Infrastructure & 
\textbf{Target 9.1:}\newline Quality, reliable infrastructure. & 
[Project contribution to resilient infrastructure design and capital efficiency.] \\
\bottomrule
\end{tabularx}}
\end{table}

% =============================================================
% TABLE OF CONTENTS, LIST OF FIGURES, LIST OF TABLES
% =============================================================
\newpage
\setstretch{1.15}
\tableofcontents

\newpage
\begingroup
\let\clearpage\relax
\listoffigures
\vspace{0.8cm}
\listoftables
\endgroup

% =============================================================
% MAIN CONTENT CHAPTERS (Arabic Numerals, Continuous Section Flow)
% =============================================================
\newpage
\pagenumbering{arabic}
\setcounter{page}{1}
\setstretch{1.25}

\chapter{Introduction}
\label{chap:intro}

\section{Background and Context}
[Provide introductory background on the domain and application.]

\section{Motivation}
[Describe the motivation for employing computational graph algorithms.]

\section{Scope of the Project}
[Delineate the explicit engineering scope.]

\section{Organization of the Report}
[Brief chapter-by-chapter outline of the report.]

\chapter{Problem Definition and Objectives}
\label{chap:problem}

\section{Problem Statement}
[State the exact mathematical/algorithmic problem statement.]

\section{Objectives}
\begin{itemize}
  \item Objective 1
  \item Objective 2
  \item Objective 3
  \item Objective 4
\end{itemize}

\chapter{Dataset and Network Preprocessing}
\label{chap:dataset}

\section{Data Acquisition}
[Describe source data extraction and preprocessing pipeline.]

\section{Graph Model Formulation}
[Define the vertex and directed edge attributes.]

\chapter{Exploratory Graph Analysis}
\label{chap:eda}

\section{Topological Properties}
[Examine degree distributions, sparsity, clustering, and connectivity.]

\chapter{Algorithmic Methodology}
\label{chap:methodology}

\section{Core Algorithms}
[Detail mathematical formulations, equilibrium conditions, and search heuristics.]

\chapter{Experimental Results and Performance Analysis}
\label{chap:results}

\section{Evaluation Metrics}
[Define quantitative performance indicators.]

\section{Comparative Experimental Results}
[Present experimental results across baseline and surge traffic scenarios.]

\chapter{System Architecture and Implementation}
\label{chap:implementation}

\section{Architecture Overview}
[Detail system components, backend APIs, and frontend GIS interface.]

\chapter{Discussion and Sensitivity Analysis}
\label{chap:discussion}

\section{Engineering Insights}
[Interpret systemic implications, physical tradeoffs, and computational scaling.]

\chapter{Limitations and Future Research}
\label{chap:limitations}

\section{Limitations}
[State theoretical or empirical boundary conditions.]

\section{Future Directions}
[Outline future extensions and research avenues.]

\chapter{Conclusion}
\label{chap:conclusion}
[Summarize project achievements, concrete findings, and engineering takeaways.]

% =============================================================
% REFERENCES
% =============================================================
\newpage
\begin{thebibliography}{99}
\addcontentsline{toc}{chapter}{References}

\bibitem{ref1}
Author Name,
\newblock ``Paper or Book Title,''
\newblock {\em Journal / Publisher}, vol.~1, no.~1, pp.~1--10, 2026.

\end{thebibliography}

\end{document}
